\chapter{Conclusiones y trabajo futuro}\label{sec:conclusion}

\section{Respuestas a las preguntas de investigación}
La tesis relacionó dos decisiones de la fusión de texturas: qué bloques combinar y para qué imágenes adquirirlos. Frente a P1, el benchmark de 22 descriptores mostró que la estrategia de mayor macro-F1 cambia con el dataset y el clasificador. La concatenación completa puede mejorar o empeorar el resultado respecto de la selección individual, mientras GFS permite reducir dimensionalidad. La selección de representaciones debe evaluarse con el clasificador que las utiliza y su eficiencia debe medirse más allá del tamaño del vector.

Frente a P2, la cascada de referencia documentó adquisición condicional sobre 27 condiciones y pilotos físicos delimitados. Las diferencias predictivas, la procedencia KTH y las discrepancias de representación CUReT/Soil limitan el alcance conjunto. Los resultados no establecen equivalencia ni una garantía de eficiencia física en todo el panel.

Para P3 se implementó la selección de base, compuerta y rescate de ensamble, reutilizando bloques compartidos y distinguiendo votos de modelos únicos. En desarrollo Outex, relevancia/B60 registró macro-F1 0,910742 frente a 0,914389 del ensamble permanente, con costo proyectado 85,40\% menor. El resultado expresa una pérdida predictiva observada a cambio de menor adquisición estimada; no acredita no inferioridad, ahorro físico o superioridad frente a Completa-22.

La aportación implementada es un procedimiento de selección bajo presupuesto medio de adquisición que distingue bloques, votos y modelos únicos. La formulación explicita cuándo una ruta condicional puede ahorrar adquisición respecto de ejecutar la misma unión permanentemente. El control del mismo rescate distingue una reducción por poda de 69,76\% y otra sucesiva por compuerta de 51,72\% en relevancia/B60, con pequeñas pérdidas predictivas observadas. En marginal/B60 se conservan las 680 predicciones con 43,09\% menos adquisición proyectada, aunque aumentan las llamadas. La evidencia separa esos mecanismos en este archivo de desarrollo; no establece ahorro físico ni generalización del híbrido a los seis conjuntos.

Frente al ensamble estático seleccionado internamente bajo el mismo presupuesto, relevancia/B60 obtiene en el piloto macro-F1 $0{,}910742$ frente a $0{,}892923$, con costos proyectados de $53{,}316$ y $44{,}831$ ms. En B15 y B120 el control estático alcanza mayor F1 con mayor costo. La adquisición condicional aporta, por tanto, puntos competitivos de desarrollo dentro de la familia examinada, no superioridad uniforme ni optimalidad global.

La réplica con tres semillas FMD y dos semillas adicionales Outex mantiene un compromiso medio favorable de relevancia/B60 frente al estático: $0{,}977489/35{,}689$ ms frente a $0{,}974600/41{,}888$ ms en FMD, y $0{,}905991/40{,}750$ ms frente a $0{,}889152/49{,}047$ ms en las nuevas semillas Outex. FMD/42 favorece al estático en ambas medidas, y relevancia tampoco supera siempre a marginal. La activación del mismo rescate mejora F1 medio en FMD/B60 y lo reduce en Outex. La réplica amplía el piloto a dos datasets, pero no convierte sus semillas en dominios independientes ni acredita el híbrido en los seis conjuntos.

Finalmente, el cierre del panel físico online sobre las 94 condiciones únicas (362.640 llamadas individuales en GPU RTX 4060) y la comparación empírica contra cuatro métodos del estado del arte basados en concatenación —\texttt{Neshov BEiT-v2} multicapa \citep{neshov2025aggregation}, \emph{DeepLumina} \citep{simon2022deeplumina}, \emph{HyTexNet} \citep{gupta2026hytexnet} y \emph{6LBP+dCellNN} \citep{chang2024cellnn}— consolidan la validación externa de la tesis. Tras verificar la paridad matemática de las implementaciones ($\text{error} = 0{,}0$ frente al código y pesos oficiales en \emph{6LBP+dCellNN}, reproducción exacta de la Tabla 1 de \emph{HyTexNet} y validación del rango publicado de $90{,}00\%$ en FMD 80/20 y $98{,}81\%$ en CUReT para \emph{DeepLumina}), la evaluación en el panel común demuestra que nuestras políticas de adquisición condicional (\texttt{Adaptive-selection-margin}) y de ensamble (\texttt{Ens-diversidad} y \texttt{Ens-tombola}) superan en macro-F1 a los cuatro comparadores externos en la totalidad de los conjuntos de datos evaluados, logrando además menor latencia física por imagen ($21{,}65$--$54{,}35$ ms) que las concatenaciones de múltiples pasadas como \emph{DeepLumina} ($51{,}57$--$57{,}80$ ms) y reducciones de tiempo de un orden de magnitud frente a la concatenación estática completa.

\section{Limitaciones y trabajo futuro}
La comparación presupuestada debe ampliarse más allá del ensamble estático de hasta tres miembros del pool visitado y de los dos datasets de desarrollo del híbrido. Se requieren otros dominios, sensibilidad a la base y controles GFS/Top-$k$ concatenados bajo el mismo protocolo. El presupuesto de despliegue y la tolerancia predictiva deben fijarse antes de evaluar nuevas condiciones, no inferirse retrospectivamente del punto más favorable. La sensibilidad sin KTH reduce una reserva del análisis histórico, pero no reconcilia su protocolo de origen.

La validación de eficiencia requiere ejecutar la ruta completa desde imagen hasta predicción, incluyendo clasificadores, transferencias y un régimen explícito de residencia. Antes de generalizar las mediciones deben resolverse o delimitarse la fidelidad CUReT/Soil y la procedencia KTH. Estas verificaciones sostienen afirmaciones centrales pendientes; no se consideran realizadas por incluirlas como trabajo futuro.

Una extensión posterior puede integrar representaciones multicapa, evaluar energía o estudiar dispositivos de borde y políticas de más etapas. Esas ampliaciones serán útiles si modifican de forma verificable el compromiso calidad--costo. La conclusión central del estudio es que optimizar la fusión requiere distinguir la representación seleccionada de su adquisición por imagen y evaluar ambas decisiones con controles comparables.
